\begin{figure}[htbp]
\centering
\begin{tikzpicture}[
    tn/.style={base, text width=1.42cm, minimum height=1.5cm, font=\scriptsize},
    over/.style={preaction={draw=white, line width=3.2pt, -}},
  ]
  \def\dx{2.04}
  \def\ya{0}\def\yb{-3.35}\def\yc{-7.15}

  % ---------- stage A ----------
  \node[grp, anchor=west] at (-0.8,{\ya+1.2}) {Stage A, still images};
  \node[tn, dat]  (a0) at (0*\dx,\ya) {still pairs:\\ LLVIP, M3FD,\\ FLIR and\\ RGBT-Tiny};
  \node[tn, n]    (a1) at (1*\dx,\ya) {random\\ frames,\\ batch 16};
  \node[tn, n]    (a2) at (2*\dx,\ya) {paired\\ geometry,\\ mosaic on};
  \node[tn, fuse] (a4) at (4*\dx,\ya) {\method\\ image mode,\\ empty state};
  \node[tn, n]    (a5) at (5*\dx,\ya) {detection\\ loss\\ $\mathcal{L}_{\text{det}}$};
  \node[tn, n]    (a6) at (6*\dx,\ya) {MuSGD,\\ AMP, EMA;\\ checkpoint A};
  \draw[arr] (a0) -- (a1); \draw[arr] (a1) -- (a2);
  \draw[arr] (a2) -- node[lbl, above=1pt] {no corruption} (a4);
  \draw[arr] (a4) -- (a5); \draw[arr] (a5) -- (a6);

  % ---------- stage B ----------
  \node[grp, anchor=west] at (-0.8,{\yb+1.2}) {Stage B, clips};
  \node[tn, dat]   (b0) at (0*\dx,\yb) {sequences:\\ RGBT-Tiny,\\ VT-MOT,\\ VisDrone VID};
  \node[tn, n]     (b1) at (1*\dx,\yb) {clips of\\ $L$ = 8 to 16,\\ stride 1--3};
  \node[tn, n]     (b2) at (2*\dx,\yb) {shared\\ geometry\\ per clip};
  \node[tn, alert] (b3) at (3*\dx,\yb) {corruption\\ simulator};
  \node[tn, tim]   (b4) at (4*\dx,\yb) {\methodS\\ unrolled,\\ burn-in 2,\\ TBPTT};
  \node[tn, n]     (b5) at (5*\dx,\yb) {loss\\ $\overline{\mathcal{L}}_{\text{det}}$ $+$\\ $\lambda_{\text{rel}}\,$BCE};
  \node[tn, n]     (b6) at (6*\dx,\yb) {optimiser;\\ FP32 scan;\\ checkpoint B};
  \foreach \i/\j in {0/1,1/2,2/3,3/4,4/5,5/6} \draw[arr] (b\i) -- (b\j);

  % simulator details above its node
  \node[lbl, anchor=south, text width=2.3cm, text=cAlert] at (3*\dx,{\yb+0.82})
       {bursty dropout, blackout,\\ drift, desync, gaps};

  % checkpoint A initialises stage B
  \draw[arr] (a6.south) -- ++(0,-0.5) -| node[lbl, pos=0.22, above=0.5pt] {initialise} ($(b4.north)+(0.35,0)$);

  % reliability targets from the simulator to the loss
  \draw[darr, draw=cAlert] (b3.south) -- ++(0,-0.42) -| (b5.south);
  \node[lbl, anchor=north, text=cAlert] at (4.55*\dx,{\yb-1.25}) {reliability targets $q$};

  % state carried to the next clip
  \draw[darr, draw=cTime, over] ($(b4.south)+(-0.35,0)$) -- ++(0,-1.25) -| (b1.south);
  \node[lbl, anchor=north, text=cTime] at (2.5*\dx,{\yb-2.08}) {final state initialises the next clip, $p=0.5$};

  % ---------- evaluation ----------
  \node[grp, anchor=west] at (-0.8,{\yc+1.15}) {Evaluation};
  \node[base, fuse, fill=cFuse!10, text width=8.4cm, minimum height=1.45cm, font=\scriptsize, align=left] (ev) at (2.45*\dx,\yc)
     {\textbf{image metrics}\enspace AP, AP\textsubscript{50}, AP\textsubscript{S}, AP below 16\,px, SAFit\\[2pt]
      \textbf{stream robustness}\enspace degradation curves, recovery time, flicker\\[2pt]
      \textbf{efficiency}\enspace latency with NMS, memory against sequence length};
  \node[tn, n, minimum height=1.45cm] (log) at (6*\dx,\yc) {run log:\\ seed, config\\ hash, tables};
  \draw[arr] (b6.south) -- ++(0,-1.6) -| ($(ev.north)+(3.4,0)$);
  \draw[arr] (ev) -- (log);
\end{tikzpicture}

\vspace{2mm}
\caption[The two-stage training pipeline]{The two-stage training pipeline.
Stage A trains the still-image \method{} and yields the Phase~1 model.
Stage B continues on clips with a corruption simulator, truncated
back-propagation through time (TBPTT) and stochastic carry-over of the state
between clips. Because the simulator knows which sensor it disabled or
shifted, it provides targets for the reliability heads. Every evaluation is
logged with its seed and configuration hash, and the tables of the final
paper are generated from these logs.}
\label{fig:training}
\end{figure}
